%%
%% Proposta de projeto -- PPC 2026/27
%% Baseado no template sample-sigchi.tex modificado pelo DI/FCUL
%%
\documentclass[sigplan]{acmart}
\settopmatter{printacmref=false} % Removes citation information below abstract
\renewcommand\footnotetextcopyrightpermission[1]{} % removes footnote with conference information in first column
\usepackage[english]{babel}
\usepackage{subcaption}
\usepackage{caption}
\usepackage{array}
\usepackage{multirow}
\usepackage{booktabs}
\usepackage{enumitem}

\begin{document}
\raggedbottom

\title{Scalability of Parallel Code in WebAssembly in the Browser vs.\ Native}

\author{David Coimbra - fc68024}
\affiliation{%
 \institution{
  Programação Paralela e Concorrente \\
  Mestrado em Engenharia Informática \\
  Faculdade de Ciências, Universidade de Lisboa}
 \country{Portugal}
 }
\email{fc68024@alunos.fc.ul.pt}

\begin{abstract}
I propose to measure whether the \emph{same} parallel code (Rust/Rayon) scales
the same way natively and in WebAssembly (WASM) with threads in the browser. An
image filter provides the baseline; a histogram varies synchronization and
contention to isolate their effect, extending the study by Kyriakou and
Tselikas~\cite{kyriakou2022}.
\end{abstract}

\keywords{WebAssembly, multithreading, scalability, synchronization, performance}

\pagestyle{plain} % removes running headers

\maketitle

\section{Introduction}

\noindent\textbf{Context and problem.} More and more heavy work runs on the web
client (images, documents, data visualization, ML). WASM with threads makes it
possible to parallelize it in the browser, but it is unknown whether it scales
like the same code natively, or whether synchronization widens the gap.

\noindent\textbf{Motivation.} As a full stack developer, I have done image
processing on the client, without threads or WASM. Is it worth parallelizing,
and for what kind of work?

\noindent\textbf{Goals.}
\begin{description}[nosep]
  \item[RQ1] How does scalability (relative speedup and efficiency) of the same
  parallel code, executed in WebAssembly in the browser and natively, evolve as
  the number of threads increases, in workloads representative of heavy
  processing on the web client?
  \item[RQ2] To what extent does synchro\-nization between threads affect the
  scalability gap between Web\-Assembly and native observed in RQ1?
\end{description}

\section{Background} \label{sec:background}

By default, Web Workers do not share memory and communicate by copying
(\texttt{postMessage}). WASM~\cite{haas2017} with threads uses a
\texttt{Shared\-Array\-Buffer} as shared memory, which requires
\emph{cross-origin isolation} (COOP/COEP); each worker instantiates the module,
and atomics only offer sequential consistency~\cite{sanchez2025}. Relative
speedup compares each environment with itself; efficiency is speedup per thread.

\section{Related Work} \label{sec:relatedwork}

In single-threaded execution, Jangda et al.~\cite{jangda2019} measure WASM to be 45--55\%
slower than native on average, and Yan et al.~\cite{yan2021} show that compiler,
flags, input and browser change the results. Outside the browser, the allocator
and atomics cause multithreading overhead~\cite{sanchez2025}. In Chromium,
Kyriakou and Tselikas~\cite{kyriakou2022} measure 93.5\% of the speed of a
Rust addon in Node.js single-threaded and 67.9\% multithreaded, but with a
single workload without synchronization, 1 vs.\ 24 threads, and means only.

\section{Experimental Design} \label{sec:work1}

\noindent The same Rust code (Rayon), with the same compiler and flags, runs
natively and in WASM on Chrome (V8) and Firefox (SpiderMonkey), with 1, 2,
4\ldots\ $N$ threads in each environment.

\noindent\textbf{RQ1 - does WASM scale like native?} I compare the speedup and
efficiency curves of all workloads. The reference is the \emph{image filter}
(embarrassingly parallel, no explicit synchronization): WASM's baseline cost in the browser.

\noindent\textbf{RQ2 - does synchronization push WASM away from native?}
Workload: \emph{image histogram}, 2$\times$2 grid:
\begin{itemize}[nosep,leftmargin=*]
  \item \emph{Main axis - with vs.\ without synchronization} (shared atomics
  vs.\ per-thread local counters, the control): same problem and same result,
  only synchronization changes. It tells \textbf{whether} it causes the gap.
  \item \emph{Second axis - low vs.\ high contention} (noise image vs.\ few
  colors): it tells \textbf{how much} contention matters.
\end{itemize}

\noindent\textbf{Controls and analysis.} Inputs with a fixed seed, checksum,
pool created before measuring, warm-up, i9-14900KF on Linux. Commercial browsers at fixed versions; P/E cores (and maximum $N$) chosen through preliminary experiments in the pilot. Independent runs in random order, statistical
tests and effect size~\cite{georges2007,hoefler2015}.

\section{Forthcoming Work and Conclusions} \label{sec:conclusions}

\noindent\textbf{Plan.} Pilot (Oct 15--25), harness and both workloads with
preliminary data (check-in, Nov 11), campaign and analysis (Nov 20), paper
(Dec 2).

\noindent\textbf{Use of AI.} I used Claude to help with the design and to review and
polish the text; the decisions are mine.

\bibliographystyle{ACM-Reference-Format}
\bibliography{../bib/referencias}

\end{document}
